\section{Finite fields and field-size scaling}
\label{sec:prime-fields}

Over $\F_2$, a nonzero message is determined by its support.  Over
$\F_p$, two messages with the same support can cancel differently.  An
addition-only expander therefore has no support-only output law.  We restore
that symmetry by assigning an independent label from $\F_p^*$ to every edge.

Fix a prime power $p$, choose the field $\F_p$, and write $s:=p-1$.
Multiplying a fixed nonzero message value
by a uniform edge label gives a uniform element of $\F_p^*$.  Thus random
labels erase the message values while preserving its support.

The analysis has two scales.  The asymptotic theorem below fixes $p$ while
the code length grows.  The finite certificates in
Section~\ref{sec:certificates} instead fix a cryptographic-size length and
fields whose orders are near $2^{128}$.  In that finite regime, the proof must
control the projective message count rather than absorb it into an
asymptotic exponent.

\subsection{Labeled Bernoulli incidence}

In the labeled Bernoulli ensemble, an entry is zero with probability
$1-\rho$ and equals each element of $\F_p^*$ with probability $\rho/s$.

\begin{lemma}[Labeled Bernoulli output law]
\label{lem:p-bernoulli-law}
Fix $\x\in\F_p^k$ with support size $r$.  The coordinates of $\x\B$ are
independent.  Each coordinate is nonzero with probability
\begin{equation}
 a^{(p)}_r:=\frac{s}{p}
 \left(1-\left(1-\frac{p\rho}{s}\right)^r\right).
 \label{eq:p-activation}
\end{equation}
Conditional on being nonzero, its value is uniform in $\F_p^*$.
\end{lemma}

\begin{proof}
Let $\chi$ be a nonconstant additive character of $\F_p$, that is, a
complex-valued function satisfying $\chi(x+y)=\chi(x)\chi(y)$.  One active row
contribution $Z$ satisfies
\[
 \E[\chi(Z)]
 =1-\rho+\frac{\rho}{s}\sum_{a\in\F_p^*}\chi(a)
 =1-\frac{p\rho}{s}.
\]
The $r$ active rows are independent.  Fourier inversion over the additive
group of $\F_p$ gives \eqref{eq:p-activation} and equal mass at every
nonzero value.  Different columns use independent samples.
\end{proof}

\subsection{Finite-field EA}

Suppose an accumulator input is zero with probability $1-a$ and is otherwise
uniform in $\F_p^*$.  Group the accumulator states into zero and nonzero.
These two classes have the exact transfer matrix
\begin{equation}
 \mathbf T^{(p)}_a(Y):=
 \begin{pmatrix}
  1-a&aY\\
  a/s&(1-a/s)Y
 \end{pmatrix}.
 \label{eq:p-ea-transfer}
\end{equation}
From a nonzero state, exactly one of the $s$ possible nonzero increments
returns to zero.  Hence this two-state reduction is exact.

All nonzero scalar multiples of one message produce codewords with the same
support.  A \emph{projective message line} is a one-dimensional subspace of
$\F_p^k$.  The number of such lines with support size $r$ is
\begin{equation}
 N_{p,k,r}:=\binom{k}{r}s^{r-1}.
 \label{eq:projective-count}
\end{equation}

\begin{theorem}[Labeled Bernoulli EA first moment]
\label{thm:p-ea-first-moment}
Sample a labeled Bernoulli expander $\B$ and set $\G:=\B\A$.  Then
\begin{equation}
 \Pr_{\B}[\Bad_L(\G)]\leq
 \sum_{r=1}^{k}N_{p,k,r}
 \sum_{h=0}^{L}[Y^h]\,
 \e_0^{\mathsf T}
 \mathbf T^{(p)}_{a^{(p)}_r}(Y)^n\one.
 \label{eq:p-ea-first-moment}
\end{equation}
\end{theorem}

\begin{proof}
Lemma~\ref{lem:p-bernoulli-law} gives the exact two-class input law for every
projective line with support size $r$.  Matrix
\eqref{eq:p-ea-transfer} gives its output-weight generating function.  Sum
over projective lines and apply the first-moment bound.  Counting one
representative per line is sufficient because scalar multiplication does not
change output support.
\end{proof}

For $a\in[0,s/p]$ and $z\in(0,1]$, define
\begin{align}
 \lambda_p(a,z)
 &:=\frac12\left(1-a+(1-a/s)z
 +\sqrt{(1-a-(1-a/s)z)^2+4a^2z/s}\right),
 \label{eq:p-ea-spectral-radius}\\
 J^{(p)}_\delta(a)
 &:=\sup_{0<z\leq1}
 \{\delta\ln z-\ln\lambda_p(a,z)\},
 \label{eq:p-ea-rate-function}\\
 I_{p,\delta}
 &:=\left(\sqrt{1-\delta}
 -\sqrt{\frac{\delta}{p-1}}\right)^2.
 \label{eq:p-ea-sparse-rate}
\end{align}

\begin{lemma}[Finite-field accumulator endpoints]
\label{lem:p-ea-endpoints}
For $0<\delta<s/p$,
\begin{align}
 \Pr[H\leq\lfloor\delta n\rfloor]
 &\leq\sqrt p\exp(-nJ^{(p)}_\delta(a)),
 \label{eq:p-ea-tail}\\
 \liminf_{a\downarrow0}\frac{J^{(p)}_\delta(a)}a
 &\geq I_{p,\delta},
 \label{eq:p-ea-sparse-endpoint}\\
 J^{(p)}_\delta(s/p)&=\ln p-h_p(\delta).
 \label{eq:p-ea-dense-endpoint}
\end{align}
Here $H$ is the output weight for independent inputs that are nonzero with
probability $a$ and uniform conditional on being nonzero.
\end{lemma}

\begin{proof}
Conjugating \eqref{eq:p-ea-transfer} by
$\diag(1,\sqrt{sz})$ gives a symmetric positive semidefinite matrix with
largest eigenvalue $\lambda_p(a,z)$.  The same positive-marker argument as in
Lemma~\ref{lem:ea-lower-tail} gives \eqref{eq:p-ea-tail}.

For the sparse endpoint, set $z=e^{-\theta a}$ and expand at $a=0$.
Optimizing the coefficient of $a$ gives
$(\sqrt{1-\delta}-\sqrt{\delta/s})^2$.  At $a=s/p$, the input and accumulator
state are uniform in $\F_p$, and
$\lambda_p(s/p,z)=(1+sz)/p$.  The optimizer
$z=\delta/(s(1-\delta))$ gives \eqref{eq:p-ea-dense-endpoint}.
\end{proof}

\begin{theorem}[Fixed-field EA distance]
\label{thm:p-ea-distance}
Fix a prime power $p$ and the field $\F_p$.  Let $k/n\to R\in(0,1)$ and sample a labeled Bernoulli
expander of density $\rho_n=C\ln(n)/n$.  If
$0<\delta<(p-1)/p$ and
\begin{align}
 C I_{p,\delta}&>1,
 \label{eq:p-ea-sparse-condition}\\
 h_p(\delta)&<(1-R)\ln p,
 \label{eq:p-ea-gv-condition}
\end{align}
then
\begin{equation}
 \Pr_{\B}[d_{\min}(\B\A)\leq\lfloor\delta n\rfloor]=o(1).
 \label{eq:p-ea-distance-conclusion}
\end{equation}
\end{theorem}

\begin{proof}
Use the three support ranges from Theorem~\ref{thm:ea-distance}.  For sparse
supports, $a^{(p)}_r=r\rho_n(1-o(1))$, while
$N_{p,k,r}\leq(sk)^r$.  Equations \eqref{eq:p-ea-sparse-endpoint} and
\eqref{eq:p-ea-sparse-condition} make their sum vanish.  For intermediate
supports, the rate function is uniformly positive and dominates their small
counting exponent.  For $r>\eta n$, one has
$a^{(p)}_r=s/p-o(1)$ uniformly.  There are at most $p^k/s$ projective lines,
so their total exponent is
\[
 R\ln p+h_p(\delta)-\ln p+o(1),
\]
which is negative by \eqref{eq:p-ea-gv-condition}.
\end{proof}

The phrase ``fix a field'' is essential.  If $p$ grows with $n$, the factor
$s^{r-1}$ in \eqref{eq:projective-count} changes the sparse counting exponent.
Theorem~\ref{thm:p-ea-distance} therefore does not justify a cryptographic-size
field at finite length.  The next bounds group projective messages more
carefully.

\subsection{Constraint traces for the projective union}

For a fixed support $S$, the elementary union bound multiplies a per-line
failure probability by $s^{|S|-1}$.  That multiplier can exceed one long
before the underlying event becomes likely.  The first correction is to cap
the union for each support:
\begin{equation}
 \Pr[\Bad_L(\B\A)]
 \leq\sum_{r=1}^{k}\binom{k}{r}
 \min\{1,s^{r-1}\tau_r(L)\},
 \label{eq:support-grouping}
\end{equation}
where $\tau_r(L)$ is the bad-output probability of one projective line.
Here RAA denotes repeat--accumulate--accumulate codes.  This support grouping
is related to recent large-field RAA
analyses~\cite{AkhianiZhang2026,Khabbazian2026}.

A sharper bound counts the independent field equations imposed by one output
trace.  Call an intermediate coordinate \emph{occupied} when at least one
edge from the message support reaches it.  In the Bernoulli ensemble, an
occupied indicator has probability
\begin{equation}
 u_r:=1-(1-\rho)^r.
\end{equation}
Let $z$ mark a nonzero accumulator output and let $v$ mark an occupied
coordinate whose new accumulator state is zero.  Define
\begin{equation}
 \mathbf K_u(z,v):=
 \begin{pmatrix}1-u+uv&uz\\uv&z\end{pmatrix}.
 \label{eq:p-ea-trace-matrix}
\end{equation}

\begin{theorem}[Bernoulli constraint-trace bound]
\label{thm:p-ea-trace}
For the labeled Bernoulli EA ensemble,
\begin{equation}
 \Pr[\Bad_L(\B\A)]\leq
 \sum_{r=1}^{k}\binom{k}{r}
 \sum_{h=0}^{L}\sum_{a=0}^{n}
 [z^hv^a]\,
 \e_0^{\mathsf T}\mathbf K_{u_r}(z,v)^n\one
 \min\{1,s^{r-1-a}\}.
 \label{eq:p-ea-trace-bound}
\end{equation}
\end{theorem}

\begin{proof}
Fix a message support $S$ of size $r$, the occupied coordinates, and one
zero--nonzero accumulator trace.  Suppose $a$ occupied coordinates end in
state zero.  Each such event gives one prefix equation in the active edge
labels.  Subtracting consecutive equations isolates a nonempty set of fresh
labels, so the $a$ equations have rank $a$.  One projective line satisfies
them with probability at most $s^{-a}$.  Union-bound the $s^{r-1}$ lines on
$S$, cap the result at one, and average the traces with
\eqref{eq:p-ea-trace-matrix}.
\end{proof}

The factor $\min\{1,s^{r-1-a}\}$ is the key cancellation.  A trace with more
independent equations than projective message degrees of freedom gains a
field factor for every surplus equation.  A trace with too few equations is
charged at most one, not the full projective count.

\subsection{Two-sided regular trace for EA}

For the two-sided regular ensemble, each right coordinate has $d_R$ incident
slots.  A support of size $r$ maps to a uniform $r$-subset of the
$k=d_R\ell$ regional slots.  Define
\begin{equation}
 \mathbf Q_0(z):=\begin{pmatrix}1&0\\0&z\end{pmatrix},
 \qquad
 \mathbf Q_1(z,v):=\begin{pmatrix}v&z\\v&z\end{pmatrix},
 \label{eq:p-ea-regular-trace-inputs}
\end{equation}
and
\begin{equation}
 \mathbf R^{\EA}_{\ell,d_R,r}(z,v):=
 \binom{k}{r}^{-1}[X^r]
 \left(\mathbf Q_0(z)+
 ((1+X)^{d_R}-1)\mathbf Q_1(z,v)\right)^\ell.
 \label{eq:p-ea-biregular-region}
\end{equation}

\begin{theorem}[Two-sided regular EA trace]
\label{thm:p-ea-biregular-trace}
For the labeled two-sided regular EA ensemble,
\begin{equation}
 \Pr[\Bad_L(\B\A)]\leq
 \sum_{r=1}^{k}\binom{k}{r}
 \sum_{h=0}^{L}\sum_{a=0}^{n}[z^hv^a]\,
 \e_0^{\mathsf T}
 \mathbf R^{\EA}_{\ell,d_R,r}(z,v)^{d_L}\one
 \min\{1,s^{r-1-a}\}.
 \label{eq:p-ea-biregular-trace}
\end{equation}
\end{theorem}

\begin{proof}
In one region, $(1+X)^{d_R}-1$ counts a nonempty set of active slots at one
right coordinate.  Coefficient extraction conditions the regional slots to
contain exactly the $r$ supported left vertices.  Matrix multiplication
preserves the accumulator state across regions.  The rank argument from
Theorem~\ref{thm:p-ea-trace} supplies the projective factor.
\end{proof}

\subsection{Finite-field EC}

We now combine labeled two-sided regular incidence with the convolution in
\eqref{eq:prime-convolution}.  The full-field feedback distribution preserves
a reduced state with $m+1$ classes.  State $j<m$ records exactly $j$ trailing
zero outputs.  State $m$ means that the entire feedback state is zero.

The next transfer counts upper bounds rather than exact transition
probabilities.  Let $z$ mark a nonzero output.  Let $v$ mark a zero event that
imposes one equation in fresh randomness.  For $b\in\{0,1\}$, where $b$
records whether the expander coordinate is occupied, define matrices
$\mathbf P_b(z,v)$ by
\begin{align}
 (\mathbf P_b)_{j,0}&=z,&
 (\mathbf P_b)_{j,j+1}&=v
 &&(0\leq j<m),
 \label{eq:p-ec-active-trace}\\
 (\mathbf P_0)_{m,m}&=1,&
 (\mathbf P_1)_{m,0}&=z,&
 (\mathbf P_1)_{m,m}&=v.
 \label{eq:p-ec-zero-trace}
\end{align}

If the convolution state is nonzero, a zero output fixes one affine equation
in the fresh vector $\boldsymbol\alpha_t$.  Its probability is $1/p\leq1/s$.
If the state is zero and the coordinate is occupied, a zero expander symbol
fixes one equation in labels used only at that coordinate.  Its probability
is at most $1/s$.  All other transition probabilities are bounded by one.

Define the regional matrix
\begin{equation}
 \mathbf R^{\EC}_{\ell,d_R,r}(z,v):=
 \binom{k}{r}^{-1}[X^r]
 \left(\mathbf P_0(z,v)+
 ((1+X)^{d_R}-1)\mathbf P_1(z,v)\right)^\ell.
 \label{eq:p-ec-biregular-region}
\end{equation}

\begin{theorem}[Two-sided regular finite-field EC trace]
\label{thm:p-ec-trace}
Independently sample a labeled two-sided regular expander $\B$ and the
finite-field convolution $\C$ from \eqref{eq:prime-convolution}.  Then
\begin{equation}
 \Pr_{\B,\C}[\Bad_L(\B\C)]\leq
 \sum_{r=1}^{k}\binom{k}{r}
 \sum_{h=0}^{L}\sum_{a=0}^{n}[z^hv^a]\,
 \e_m^{\mathsf T}
 \mathbf R^{\EC}_{\ell,d_R,r}(z,v)^{d_L}\one
 \min\{1,s^{r-1-a}\}.
 \label{eq:p-ec-trace}
\end{equation}
\end{theorem}

\begin{proof}
Fix a support, one projective message line, and one output trace.  Expose the
randomness in time order.  Every $v$-marked event imposes an equation in a
fresh feedback vector or in fresh coordinate labels.  A trace with $a$ marks
therefore has probability at most $s^{-a}$ for the fixed line.  Union-bound
the $s^{r-1}$ lines on the support and cap the result at one.  The regional
coefficient in \eqref{eq:p-ec-biregular-region} averages the two-sided regular
incidence.  Matrix multiplication retains the convolution state across
regions.
\end{proof}

Sampling feedback coefficients only from $\F_p^*$ would invalidate this
reduction.  The distribution of an inner product with a nonzero state would
then depend on the support and values of that state.  Full-field sampling
makes each nontrivial inner product uniform in $\F_p$.

\paragraph{Dependence on the field.}
The trace bound uses the field only through its order $p$.  The projective
message count contributes $(p-1)^{r-1}$.  Each independent nonzero-label
equation contributes at most $(p-1)^{-1}$, and each full-field feedback
equation contributes $p^{-1}\leq(p-1)^{-1}$.  The remaining change is the
$p$-ary GV cutoff.  Consequently, the same argument applies in prime fields
and extension fields.  It does not apply to a binary matrix that is merely
reinterpreted over an extension field, because that matrix lacks the sampled
field labels and feedback coefficients.

\subsection{Singleton refinement}

The trace in Theorem~\ref{thm:p-ec-trace} treats every occupied right
coordinate alike.  In fact, a coordinate incident to exactly one supported
message symbol cannot vanish: both that symbol and its edge label are
nonzero.

Let $\mathbf P_{\mathrm{sg}}$ equal $\mathbf P_1$ except that
$(\mathbf P_{\mathrm{sg}})_{m,m}:=0$.  Replace
\eqref{eq:p-ec-biregular-region} by
\begin{equation}
 \widehat{\mathbf R}^{\EC}_{\ell,d_R,r}(z,v):=
 \binom{k}{r}^{-1}[X^r]
 \left(
 \mathbf P_0+d_RX\mathbf P_{\mathrm{sg}}+
 ((1+X)^{d_R}-1-d_RX)\mathbf P_1
 \right)^\ell.
 \label{eq:p-ec-singleton-region}
\end{equation}
Theorem~\ref{thm:p-ec-trace} remains valid with this refined matrix.

For the lowest-equation traces, one more combinatorial bound removes a
positive-marker loss.  Let $U_r$ be the number of occupied right coordinates
when a uniform $r$-subset of the $d_R\ell$ regional slots is selected.  Put
\begin{align}
 \beta_r&:=\Pr[U_r\leq\lfloor r/2\rfloor],
 \label{eq:singleton-free-probability}\\
 J_r&:=1+\left\lfloor\frac{r-1}{m}\right\rfloor,
 \qquad
 b_r:=\left\lceil\frac{n-L-(r-1)}{\ell}\right\rceil-2J_r.
 \label{eq:required-singleton-free-regions}
\end{align}

\begin{lemma}[Singleton-free regions]
\label{lem:singleton-free-regions}
Fix a support size $r$ and restrict to traces with at most $r-1$ fresh
equations.  If $b_r>0$, their contribution to the probability in
Theorem~\ref{thm:p-ec-trace} is at most
\begin{equation}
 \binom{k}{r}\binom{d_L}{b_r}\beta_r^{b_r}.
 \label{eq:singleton-free-bound}
\end{equation}
\end{lemma}

\begin{proof}
Every later entry into the all-zero convolution state requires $m$
consecutive zero outputs from a nonzero state.  Those outputs impose $m$
fresh equations.  A trace with at most $r-1$ equations therefore has at most
$J_r$ maximal all-zero-state intervals.

Such an interval can meet at most two regions containing a singleton
coordinate.  Otherwise it encounters a singleton inside a fully traversed
region and must leave the all-zero state.  If $B$ regions are singleton-free,
the all-zero-state intervals contain at most $(2J_r+B)\ell$ positions.  All
remaining zero outputs impose fresh equations.  Thus
\[
 n-L\leq(2J_r+B)\ell+r-1,
\]
which implies $B\geq b_r$.  Regional permutations are independent, and a
singleton-free region has at most $\lfloor r/2\rfloor$ occupied coordinates.
A union bound over the required regions gives \eqref{eq:singleton-free-bound}.
\end{proof}

This lemma is needed for the degree-$22$, memory-$12$ certificate.  It removes
slack in the structural branch; it does not assert that the ensemble fails
when the bound is inconclusive.
